\section{Tessellations of the euclidean
plane}\label{tessellations-of-the-euclidean-plane}

\subsection{Uniform Tilings}\label{uniform-tilings}

Uniform tilings are a larger

\begin{definition}[Vertex transitiv]

A tiling is said to be vertex transitiv if all of its vertices are
indistiguishable from each other

\end{definition}

\begin{itemize}
\tightlist
\item[$\square$]
  (playful) take some regular polygons, find arrangements to cover the
  plane without gapsing
\item[$\square$]
  what is a tesselation
\item[$\square$]
  what makes it regular
\end{itemize}

\subsection{Tilings of the plane}\label{tilings-of-the-plane}

\begin{itemize}
\tightlist
\item[$\square$]
  definition face, edge, vertex transitivity
\end{itemize}

\subsection{Regular Tesselations}\label{regular-tesselations}

\begin{itemize}
\tightlist
\item[$\square$]
  platonic tilings
\item[$\square$]
  proof that there are only three
\item[$\square$]
  Schläfli Symbol \{p, q\} describes a tiling where q p-gons meet at
  each vertex
\end{itemize}

\subsection{Semi- and Quasi-Regular
Tessellations}\label{semi--and-quasi-regular-tessellations}

\begin{itemize}
\tightlist
\item[$\square$]
  archimedian tilings

  \begin{itemize}
  \tightlist
  \item[$\square$]
    most are semi-regular derived by truncation
  \item[$\square$]
    3.6.3.6 is the only quasi-regular archimedian tiling
  \end{itemize}
\item[$\square$]
  in a later chapter we can derive semi-regular tilings by truncating,
  and quasi-regular tilings by rectifiying regular tilings.
\end{itemize}

https://en.wikipedia.org/wiki/Euclidean\_tilings\_by\_convex\_regular\_polygons\#Archimedean,\_uniform\_or\_semiregular\_tilings

\subsection{Quasi-Regular
Tessellations}\label{quasi-regular-tessellations}

\begin{itemize}
\tightlist
\item[$\square$]
  in a later chapter we can derive quasi regular tilings by rectifying
\end{itemize}

\subsection{Tessellation with Möbius
Triangles}\label{tessellation-with-muxf6bius-triangles}

\begin{itemize}
\tightlist
\item[$\square$]
  requires discussion about angle sum
\end{itemize}

A triangle tiling is a tiling made from Schwarz triangles.

A \emph{Möbius triangle} is a triangle defined by the three tuple
\((p q r) \in \mathbb{Z}^3\), which represent it's internal angles. The
angle at each vertex are given by

\[
\begin{matrix}
\alpha = \frac{\pi}{p} & \beta = \frac{\pi}{q} & \gamma = \frac{\pi}{r}
\end{matrix}
\]

Not all choices for \((p q r)\) yield a valid triangle. Must must ensure
that the angles form a valid internal angle sum. For a triangle in the
Euclidean Plane the angles must add up to \(\pi\), hence (p, q, r)
define a Euclidean triangle exactly when

\[
\alpha + \beta + \gamma = \pi \Rightarrow
\frac{1}{p} + \frac{1}{r} + \frac{1}{q} = 1
\]

We can tessellate a surface by continuously reflecting a Möbius triangle
along it's edges.

Note that all permutations of \(p, q, r\) will produce a similar
triangle and henceforth congruent tilings,
\((p, q, r) \cong (q, r, p) \cong (r, p, q)\), etc.

\subsubsection{Proof of Existence}\label{proof-of-existence}

Let \(p, q, r \in \mathbb{Z^{+}}\) be the coefficients of a möbius
triangle. We shall proof for which choices of \(p, q, r\) there exists a
valid \textbf{Eulcidean} tiling. In order to formulate an exhaustive
proof, we order the coefficients such that \(p \le q \le r\). This way
we will cover all unique tilings without repetition.

Our strategy in this proof is to consequtively limit the number of
choices for each coefficient and make a case by case analysis for each
one.

In Euclidean geomerty the angle sum of a triangle is \(\pi\), therefore
\(\frac{1}{p} + \frac{1}{q} + \frac{1}{r} = 1\) \{\#eq:test\}.

We want \(p\) to assume the its lowest possible values, i.e.~to
represent the largest possible angle in our triangle. As a lower bound
we have \(2 \le p\), since if \(p\) was one, the other coefficients
would have to be 0. Our upper bound is reached when \(p = q = r\),
therefore if we apply (\textbf{eq:test?}) we get \(p \le 3\).

We have now have found our values for \(p \in \{ 2, 3 \}\).

\textbf{Note}

\subsection{Euler Characteristic}\label{euler-characteristic}

The euler characteristic of a cell division is given by {[}Weeks,
pp.171{]} Where \(V\) is the number of vertices, \(E\) the number of
edges and \(F\) the number of faces

\(\chi = V - E + F\)

For polyhedra \(\chi=2\) since they can be embedded on the sphere
\(\mathbb{S}^2\) {[}Coxeter, pp.~10{]}

The gauss-bonnet formula gives the area of different surfaces

\(A_\text{sphere} = 2\pi\chi\)
